\slajdid{03-s009}
\begin{frame}[label=03-s009]{Piramidalna struktura dekodera}
\begin{columns}[c,onlytextwidth]
\begin{column}{.46\textwidth}\centering
% element: 03-s009:e01
\begin{tikzpicture}[x=.9cm,y=1cm,font=\fontsize{9}{11}\selectfont]
\draw[gray!60,densely dashed] (-.15,-.38) rectangle (4.95,5.70);
\node[align=center] at (.80,5.05) {Dekoder\\4/16};
\draw[DEvod] (0,2.13) rectangle (1.65,3.38);
\node at (.48,3.15) {2/4};
\foreach \n/\yy in {3/3.15,2/2.85,1/2.55,0/2.25}{
 \node[anchor=east,inner sep=1pt] at (1.65,\yy) {Y\n};
 \coordinate (r\n) at (1.65,\yy);
}
\foreach \lab/\pin/\yy in {S3/S1/2.85,S2/S0/2.55,CS/CS/2.25}{
 \draw[DEvod] (-.55,\yy) node[left,inner sep=1pt] {\lab} --(0,\yy);
 \node[anchor=west,inner sep=1pt] at (0,\yy) {\pin};
}
\foreach \g/\y in {3/4.26,2/2.84,1/1.42,0/0}{
 \draw[DEvod] (3,\y) rectangle (4.65,\y+1.25);
 \node at (3.48,\y+1.02) {2/4};
 \foreach \k/\dy in {3/1.02,2/.72,1/.42,0/.12}{
  \pgfmathtruncatemacro{\globalindex}{4*\g+\k}
  \node[anchor=east,inner sep=1pt] at (4.65,\y+\dy) {Y\k};
  \draw[DEvod] (4.65,\y+\dy)--(5.08,\y+\dy)
   node[right,inner sep=1pt] {Y\globalindex};
 }
 \foreach \pin/\dy in {S1/.72,S0/.42,CS/.12}{
  \node[anchor=west,inner sep=1pt] at (3,\y+\dy) {\pin};
 }
 \coordinate (cs\g) at (3,\y+.12);
 \draw[DEvod] (3,\y+.72)--(2.30,\y+.72);
 \draw[DEvod] (3,\y+.42)--(2.65,\y+.42);
 \ifnum\g<3
  \node[junction] at (2.30,\y+.72) {};
  \node[junction] at (2.65,\y+.42) {};
 \fi
}
\draw[DEvod] (r3)--(1.85,3.15)|-(cs3);
\draw[DEvod] (r2)--(2.00,2.85)|-(cs2);
\draw[DEvod] (r1)--(1.85,2.55)|-(cs1);
\draw[DEvod] (r0)--(2.00,2.25)|-(cs0);
\draw[DEvod] (2.30,4.98)--(2.30,-.15)--(-.55,-.15) node[left,inner sep=1pt] {S1};
\draw[DEvod] (2.65,4.68)--(2.65,-.45)--(-.55,-.45) node[left,inner sep=1pt] {S0};
\end{tikzpicture}

\end{column}
\begin{column}{.51\textwidth}
% element: 03-s009:e02
Za povezivanje je dovoljan \textbf{jedan CS ulaz po komponenti}.
\par\vspace{1mm}
U izlaznom nivou koriste se signali \textbf{najmanje težine}, a u prethodnim nivoima
signali sve veće težine. Takav raspored olakšava određivanje indeksa novih izlaza;
mogući su i drugi načini povezivanja.
\par\vspace{1mm}
% element: 03-s009:e03
\Rezultat{Veći broj izlaza dobijamo dodavanjem nivoa.
Dekoder \(6/64\), realizovan dekoderima 2/4, ima \textbf{tri nivoa}.}
\par\vspace{1mm}
Naziv \emph{piramidalna} potiče od oblika grananja i načina dekodovanja.
\end{column}
\end{columns}
\end{frame}
